% ==============================================================================
% MÓDULO CAPÍTULO II: DESCRIPCIÓN GEOGRÁFICA, TERRITORIAL Y AMBIENTAL
% Archivo: cap_2/2.1_geografia_territorio.tex
% Enfoque: 100% Calidad de Atención (Donabedian y Watson)
% ==============================================================================
\subsection{Descripción geográfica, territorial y ambiental del área de estudio}
\label{sec:2.1_geografia}

La investigación se sitúa en el Centro de Salud Belén, establecimiento sanitario cabecera perteneciente a la Microred Huamanga de la Red de Salud Huamanga, adscrito a la Dirección Regional de Salud (DIRESA) de Ayacucho. El establecimiento se localiza en el emblemático e histórico barrio de Belén, en el sector céntrico-sur del distrito de Ayacucho, provincia de Huamanga, departamento de Ayacucho, a una altitud oficial de 2,760 metros sobre el nivel del mar. Sus coordenadas geográficas de emplazamiento corresponden a los 13$^\circ$09'47'' de Latitud Sur y 74$^\circ$13'28'' de Longitud Oeste.

El territorio bajo su responsabilidad sanitaria se asienta en la subcuenca del río Alameda, dentro de la cuenca principal del valle de Huamanga. Limita por el norte con el casco monumental y urbano central de la ciudad de Huamanga; por el sur con las urbanizaciones periurbanas del distrito de San Juan Bautista; por el este con las faldas del cerro Acuchimay y el distrito de Carmen Alto; y por el oeste con la prolongación de la avenida Mariscal Cáceres y el jirón Bellido. Su orografía presenta una topografía moderadamente accidentada, con calles estrechas, empinadas y trazados coloniales que imponen un esfuerzo físico considerable a los usuarios con movilidad reducida, gestantes y personas de la tercera edad que acuden a pie.

El clima del área es templado y seco durante las horas diurnas soleadas, con temperaturas promedio anuales que oscilan entre los 15~°C y 22~°C. No obstante, una característica meteorológica decisiva para la experiencia del usuario radica en la marcada oscilación térmica interdiurna: durante las madrugadas andinas (entre las 4:00~a.m. y las 6:30~a.m.), el termómetro desciende habitualmente a valores de entre 5~°C y 8~°C, llegando a extremos menores en la temporada seca y de heladas (junio a agosto). Asimismo, se presenta una temporada de precipitaciones pluviales intensas entre los meses de noviembre y marzo. Esta condición ambiental no constituye un dato accesorio, sino un determinante corpóreo directo que condiciona la expectativa del cuidado: el usuario que arriba al establecimiento con frío penetrante, fatiga somática y dolor acumulado demanda una acogida asistencial empática y cálida que alivie su vulnerabilidad.

En cuanto a la conectividad y movilidad urbana, el establecimiento cuenta con acceso a través de vías pavimentadas (Jr.~Bellido, Jr.~Arequipa y Av.~Mariscal Cáceres) y rutas de transporte público (líneas 1, 3, 7 y 12, además de mototaxis). A pesar de ello, las limitaciones económicas obligan a numerosas familias de sectores altos a realizar traslados peatonales de 20 a 40 minutos cargando niños en mantas tradicionales (\textit{llicllas}). Esta exigencia física y ambiental predispone un estado de agotamiento al momento de ingresar a los consultorios, transformando la calidez del trato, la prontitud del saludo y la dignidad del espacio de espera en elementos determinantes de la calidad del cuidado.
